% Options for packages loaded elsewhere
\PassOptionsToPackage{unicode}{hyperref}
\PassOptionsToPackage{hyphens}{url}
\documentclass[
]{article}
\usepackage{xcolor}
\usepackage{amsmath,amssymb}
\setcounter{secnumdepth}{-\maxdimen} % remove section numbering
\usepackage{iftex}
\ifPDFTeX
  \usepackage[T1]{fontenc}
  \usepackage[utf8]{inputenc}
  \usepackage{textcomp} % provide euro and other symbols
\else % if luatex or xetex
  \usepackage{unicode-math} % this also loads fontspec
  \defaultfontfeatures{Scale=MatchLowercase}
  \defaultfontfeatures[\rmfamily]{Ligatures=TeX,Scale=1}
\fi
\usepackage{lmodern}
\ifPDFTeX\else
  % xetex/luatex font selection
\fi
% Use upquote if available, for straight quotes in verbatim environments
\IfFileExists{upquote.sty}{\usepackage{upquote}}{}
\IfFileExists{microtype.sty}{% use microtype if available
  \usepackage[]{microtype}
  \UseMicrotypeSet[protrusion]{basicmath} % disable protrusion for tt fonts
}{}
\makeatletter
\@ifundefined{KOMAClassName}{% if non-KOMA class
  \IfFileExists{parskip.sty}{%
    \usepackage{parskip}
  }{% else
    \setlength{\parindent}{0pt}
    \setlength{\parskip}{6pt plus 2pt minus 1pt}}
}{% if KOMA class
  \KOMAoptions{parskip=half}}
\makeatother
\usepackage{longtable,booktabs,array}
\newcounter{none} % for unnumbered tables
\usepackage{calc} % for calculating minipage widths
% Correct order of tables after \paragraph or \subparagraph
\usepackage{etoolbox}
\makeatletter
\patchcmd\longtable{\par}{\if@noskipsec\mbox{}\fi\par}{}{}
\makeatother
% Allow footnotes in longtable head/foot
\IfFileExists{footnotehyper.sty}{\usepackage{footnotehyper}}{\usepackage{footnote}}
\makesavenoteenv{longtable}
\usepackage{graphicx}
\makeatletter
\newsavebox\pandoc@box
\newcommand*\pandocbounded[1]{% scales image to fit in text height/width
  \sbox\pandoc@box{#1}%
  \Gscale@div\@tempa{\textheight}{\dimexpr\ht\pandoc@box+\dp\pandoc@box\relax}%
  \Gscale@div\@tempb{\linewidth}{\wd\pandoc@box}%
  \ifdim\@tempb\p@<\@tempa\p@\let\@tempa\@tempb\fi% select the smaller of both
  \ifdim\@tempa\p@<\p@\scalebox{\@tempa}{\usebox\pandoc@box}%
  \else\usebox{\pandoc@box}%
  \fi%
}
% Set default figure placement to htbp
\def\fps@figure{htbp}
\makeatother
\usepackage{svg}
\setlength{\emergencystretch}{3em} % prevent overfull lines
\providecommand{\tightlist}{%
  \setlength{\itemsep}{0pt}\setlength{\parskip}{0pt}}
\usepackage{bookmark}
\IfFileExists{xurl.sty}{\usepackage{xurl}}{} % add URL line breaks if available
\urlstyle{same}
\hypersetup{
  hidelinks,
  pdfcreator={LaTeX via pandoc}}

\author{}
\date{}

\begin{document}

\section{OntoBridg: Fail-Closed Evaluation of Open-Relation Ontology
Learning}\label{ontobridg-fail-closed-evaluation-of-open-relation-ontology-learning}

\textbf{Keywords:} ontology learning; large language models;
neuro-symbolic systems; relation extraction; calibration; benchmark
contamination

\subsection{Abstract}\label{abstract}

Large language models can generate ontology triples from text, but
open-relation generation creates a difficult precision problem: a model
may emit many plausible, evidence-linked relations that do not match the
benchmark graph. Public ontology resources also create plausible
pretraining-exposure confounds. We present \textbf{OntoBridg}, an
ontology-learning pipeline that treats the language model as a candidate
generator and applies relation-channel evidence, confidence,
canonicalization, self-loop, predicate, and quantity constraints. We
evaluate it on the public LLMs4OL 2026 Task A and Task B training data
under a gold-withheld prompt protocol; no result uses hidden test labels
or constitutes a leaderboard claim. On relation-density-stratified
samples, isolating constraints to the relation channel improves exact
Graph Similarity by 0.0315 and 0.0409 for Tasks A and B, but exact
non-taxonomic F1 remains zero. On a frozen synthetic nonce holdout,
guarding raises non-taxonomic F1 from 0.2612 to 0.8571 with no errors on
60 negated controls. This gain does not transfer to a disjoint,
relation-rich official-text slice, where both systems recover 0 of 101
gold relations. Predicate routing recovers two \texttt{type} edges but
produces 138 official exact false positives and fails a precommitted
five-sample support gate. A contamination pilot finds high factual
probeability but opposing alias effects across tasks. These results
support a fail-closed evaluation protocol rather than a naturalistic
performance claim: controlled extraction succeeds, official-text
relation recovery remains unsupported, and every calibration and
deployment gate stays closed.

\subsection{1 Introduction}\label{introduction}

Ontology learning converts text into typed entities, taxonomic edges,
and non-taxonomic relations. The LLMs4OL program established common
tasks and datasets for testing whether large language models (LLMs) can
perform these operations across lexical, geographic, biomedical, and
web-schema domains {[}1,2{]}. More recent systems combine prompting with
retrieval or learned graph components. Beliaeva and Rahmatullaev use
retrieval-augmented prompting, embedding ensembles, and cross-attention
taxonomy induction {[}3{]}. Mishra et al.~report strong 2026 Task A and
Task B scores with retrieval-augmented Qwen2.5-14B, yet extract no
non-taxonomic relations {[}4{]}. This exposes a central design tension:
closed vocabularies can suppress false positives, while open
vocabularies are needed when the relevant property inventory is
incomplete.

Open generation alone does not solve the problem. A model can extract a
relation that is linguistically plausible, grounded in the supplied
passage, and still fail exact graph scoring because its predicate,
direction, endpoint form, or relation selection differs from the
reference graph. In ontology learning, indiscriminate recall is
especially costly because every extra edge changes graph neighborhoods
and consistency. Evidence constraints can remove unsupported candidates,
but they do not determine which of several text-supported facts belongs
in a sparse benchmark graph.

Evaluation is complicated further by possible training-data exposure.
WordNet, GeoNames, schema.org, and many biomedical ontologies are
public, widely mirrored resources. Closed-book recall of their facts is
therefore not proof of induction from the supplied text. Black-box
contamination tests can establish suspicious distributional evidence
without proving that a specific item occurred in a proprietary training
corpus {[}10,11{]}. Task-level contamination can also make nominally
few-shot evaluation behave like recall {[}12{]}. For ontology learning,
this distinction matters because the benchmark target is itself a public
knowledge artifact.

We study these issues through \textbf{OntoBridg}, a pipeline in which
the LLM generates candidate triples and deterministic postprocessing
governs whether relation candidates survive. Our goal is not to claim a
new leaderboard score. Instead, we ask where relation constraints work,
where they fail, and whether a manifest-frozen fail-closed protocol
prevents a controlled positive result from being promoted into an
unsupported deployment claim.

The paper makes four contributions:

\begin{enumerate}
\def\labelenumi{\arabic{enumi}.}
\item
  \textbf{A gold-withheld failure characterization.} On fixed prefixes
  of the public 2026 Task A and Task B training sets, gold is excluded
  from prompts and unconstrained open-relation generation emits 384 and
  150 non-taxonomic predictions. It misses Task A's sole non-taxonomic
  gold edge, while the Task B prefix contains no such gold edge; exact
  Graph Similarity falls sharply in both tasks.
\item
  \textbf{A relation-channel guard.} Confidence, verbatim evidence,
  entity and predicate canonicalization, self-loop removal,
  generic-predicate filtering, and a quantity cap are applied only to
  non-taxonomic relations. Preserving the typing and taxonomy channels
  isolates the effect of relation filtering and improves exact Graph
  Similarity on density-stratified samples.
\item
  \textbf{Independent controlled and transfer tests.} A frozen nonce
  holdout shows that the guarded extractor can recover new synthetic
  facts under familiar predicates. A disjoint official-text transfer
  slice and an independent predicate-routing slice show that this gain
  does not establish official relation recovery.
\item
  \textbf{A fail-closed evidence protocol.} Calibration, contamination
  diagnostics, frozen manifests, sample-grouped uncertainty estimates,
  and minimum positive-sample gates are recorded explicitly.
  Insufficient support closes the mechanism, confirmatory, and
  deployment gates instead of triggering same-slice prompt repair.
\end{enumerate}

\subsection{2 Background and Related
Work}\label{background-and-related-work}

\subsubsection{2.1 LLMs for ontology
learning}\label{llms-for-ontology-learning}

LLMs4OL formalized term typing, taxonomy discovery, and non-taxonomic
relation extraction as related ontology-learning capabilities {[}1{]}.
The 2024 challenge overview established shared-task datasets and
reporting conventions {[}2{]}, and the OntoLearner library later
provided a modular integration surface for benchmarking
ontology-learning systems {[}5{]}. The 2025 Alexbek system demonstrates
that no single architecture dominates all stages: retrieval-augmented
prompting is used for document-to-term/type extraction, embedding
ensembles for typing, and a cross-attention model for taxonomy induction
{[}3{]}. The 2026 pro-team system similarly combines retrieval with
deterministic filtering {[}4{]}. Its strong graph scores and empty
non-taxonomic relation output motivate our focus on relation-channel
behavior.

Classical ontology learning used lexical and lexico-syntactic
regularities, including Hearst patterns for hyponymy {[}6{]}. Modern LLM
systems replace fixed patterns with broader candidate generation, but
the underlying need for structural control remains. Ontology matching is
a neighboring, not identical, task; systems such as LogMap {[}7{]},
AgreementMakerLight {[}8{]}, and BERTMap {[}9{]} illustrate how lexical
or neural candidates can be combined with logical and structural
validation. OntoBridg adopts the same high-level separation between
candidate generation and deterministic acceptance while evaluating
ontology construction rather than cross-ontology correspondence.

\subsubsection{2.2 Contamination-aware
evaluation}\label{contamination-aware-evaluation}

Test-set contamination cannot generally be verified for a closed
commercial or hosted model. Oren et al.~show that aggregate ordering
statistics can provide black-box evidence of test-set exposure {[}10{]}.
Golchin and Surdeanu use guided prompting to trace contamination signals
across datasets and model families {[}11{]}. Li and Flanigan demonstrate
task contamination, where exposure to task-specific data changes the
interpretation of few-shot performance {[}12{]}. These methods motivate
diagnostics, not binary declarations that individual samples are clean
or contaminated.

Our protocol therefore separates three claims. \textbf{Probeability}
means a model discriminates official facts from constructed controls
without seeing task context. \textbf{Surface sensitivity} means
performance changes after endpoint aliases replace original names.
\textbf{Verified membership} would require direct knowledge of the
model's training corpus, which is unavailable here. We report the first
two and do not claim the third.

\subsection{3 Method}\label{method}

\subsubsection{3.1 Task interface}\label{task-interface}

We use the public LLMs4OL 2026 definitions {[}13{]}. In Task A
(Flagship), the model receives a text context and predicts primitive
ontology triples. In Task B (Reuse), it receives a context, an initial
ontology, and the officially supplied term and type lists, and predicts
extension triples. Let (x\_i) denote the legal input fields for sample
(i), (G\_i) its locally held gold graph, and (hat\{G\}\_i) the predicted
graph. Gold triples are used for selection or scoring only and never
enter prompts.

Triples are partitioned into typing, taxonomy, and non-taxonomic
channels. The relation experiments exclude self-loops and treat the
exact normalized triple

{[} (s,p,o) \in \hat{G}\_i \cap G\_i {]}

as a true positive. Surface or semantic variants are not counted by the
primary exact metric.

\subsubsection{3.2 Candidate-generation
systems}\label{candidate-generation-systems}

The experiments compare the following system identities:

\begin{itemize}
\tightlist
\item
  \textbf{Closed vocabulary:} generate typing and taxonomy triples only.
\item
  \textbf{Open vocabulary:} generate typing, taxonomy, and concise
  non-taxonomic predicates supported by the context.
\item
  \textbf{Guarded relation generation:} require every relation candidate
  to include a confidence score and verbatim evidence.
\item
  \textbf{Relation-channel hybrid:} retain typing and taxonomy triples
  from the legacy open system and substitute only the guarded relation
  channel.
\item
  \textbf{Predicate-routed relation generation:} restrict a
  relation-only prompt to a frozen eight-predicate inventory for an
  independent mechanism test.
\end{itemize}

The hosted backbone is identified by the stable experiment alias
\texttt{dots3-note-prev}, with temperature zero. Prompt hashes,
model-input hashes, postprocessing signatures, attempt counts, and cache
identities are retained in the experiment artifacts. Because a hosted
alias does not guarantee immutable provider weights, model drift remains
a limitation.

\subsubsection{3.3 Relation guard}\label{relation-guard}

Each guarded candidate contains a triple (r=(s,p,o)), a raw confidence
(q\in[0,1]), an evidence mode, and one or two verbatim evidence spans.
Entity strings are normalized against exact context surfaces and the
official input inventory. Predicate aliases are canonicalized with
direction-aware reversal where required.

An explicit candidate requires one verbatim span containing both
endpoints and raw confidence (q\geq0.85). A cross-sentence candidate
requires two distinct verbatim spans, joint endpoint coverage, a maximum
normalized gap of 600 characters, and (q\geq0.95). Candidates are
rejected for invalid or non-verbatim evidence, missing endpoints,
self-loops, generic predicates, or failure of an optional frozen
predicate policy. Surviving candidates are ranked by confidence and
capped at eight relations per sample. Formally,

{[} \operatorname{accept}(r)=I\_q I\_e I\_\{\neg self\}
I\_\{\neg generic\} I\_p I\_k, {]}

where (I\_q) is the mode-specific confidence test, (I\_e) is the
evidence test, (I\_p) is the optional predicate policy, and (I\_k)
represents the per-sample quantity cap. The guard is deterministic once
the raw generation is cached.

Applying the guard directly to a full generated graph can also change
typing and taxonomy behavior because the prompt itself changes. The
hybrid comparison avoids this confound: it holds the legacy
typing/taxonomy output fixed and changes only accepted relations.

\subsubsection{3.4 Calibration and fail-closed
gates}\label{calibration-and-fail-closed-gates}

Raw generation and postprocessing have separate cache identities. A new
calibrator can therefore be fit and replayed locally without another
model call. Calibration is considered deployable only when its labels
are independent of the evaluation slice and exact positives occur in at
least five independent samples. Samples, rather than candidates, define
cross-validation and bootstrap clusters.

Passing a probability metric is not sufficient. A confirmatory gate also
requires that a frozen decision rule retain exact relations in at least
five independent samples on an untouched slice. Resubstitution,
same-slice threshold search, weak labels without human audit, and
post-hoc surface rules cannot open the gate. Every artifact records
\texttt{scientific\_deployable} and \texttt{confirmatory\_eligible};
both default to false.

\subsubsection{3.5 Contamination
diagnostics}\label{contamination-diagnostics}

The closed-book probe asks the model to discriminate official facts from
matched constructed controls without task context. We summarize
discrimination with AUROC and a sample-cluster bootstrap interval. A
paired surface test replaces endpoint names with aliases while
preserving graph structure. Neither test provides verified training-set
membership. Probe results are descriptive and cannot tune generation,
filtering, calibration, or deployment.

\subsection{4 Experimental Design}\label{experimental-design}

\subsubsection{4.1 Data and frozen evidence
partitions}\label{data-and-frozen-evidence-partitions}

All official-data experiments use the public LLMs4OL 2026 Task A and
Task B training JSON files {[}13{]}. The challenge repository is
MIT-licensed; upstream ontology resources, especially UMLS-derived
material, may impose additional terms. We therefore release scripts,
manifests, hashes, prompts, and aggregate results, but do not
redistribute third-party source data in the artifact.

Table 1 separates the evidence partitions. Gold-stratified sets are
useful diagnostics but are not unbiased estimates of the natural task
distribution.

\textbf{Table 1. Frozen evidence partitions.}

{\def\LTcaptype{none} % do not increment counter
\begin{longtable}[]{@{}
  >{\raggedright\arraybackslash}p{(\linewidth - 8\tabcolsep) * \real{0.1875}}
  >{\raggedright\arraybackslash}p{(\linewidth - 8\tabcolsep) * \real{0.1875}}
  >{\raggedleft\arraybackslash}p{(\linewidth - 8\tabcolsep) * \real{0.2500}}
  >{\raggedright\arraybackslash}p{(\linewidth - 8\tabcolsep) * \real{0.1875}}
  >{\raggedright\arraybackslash}p{(\linewidth - 8\tabcolsep) * \real{0.1875}}@{}}
\toprule\noalign{}
\begin{minipage}[b]{\linewidth}\raggedright
Partition
\end{minipage} & \begin{minipage}[b]{\linewidth}\raggedright
Source
\end{minipage} & \begin{minipage}[b]{\linewidth}\raggedleft
Samples
\end{minipage} & \begin{minipage}[b]{\linewidth}\raggedright
Purpose
\end{minipage} & \begin{minipage}[b]{\linewidth}\raggedright
Selection status
\end{minipage} \\
\midrule\noalign{}
\endhead
\bottomrule\noalign{}
\endlastfoot
Blind prefix & Task A + B train & 30 per task & Closed vs.~open baseline
& Fixed prefix; gold withheld from prompts \\
Density strata & Task A + B train & 15 per task & Relation-only guard
comparison & Gold-stratified low/medium/high density \\
Independent v2 calibration & Task A + B train & 39 & Candidate
calibration support & Disjoint; relation-enriched \\
V5 confirmatory & Task A train & 24 & Frozen decision-rule test &
Disjoint; predicate-balanced, endpoint-visible \\
Nonce holdout & Synthetic context-only & 30 & Controlled new-fact
extraction & Frozen before generation \\
Official transfer & Task A train & 30 & Transfer of nonce result &
Disjoint; relation-rich, gold-stratified \\
Predicate routing & Task A train & 27 & Independent routing mechanism
test & Disjoint 8/8/3/8 predicate strata \\
\end{longtable}
}

The nonce holdout contains 120 balanced positive relations, 60 negated
co-occurrence controls, and 300 newly generated opaque entities. It has
zero exact entity-string and gold-triple overlap with the official Task
A/B training data. It is synthetic and low-overlap, not a verified
temporal benchmark.

\subsubsection{4.2 Metrics and
uncertainty}\label{metrics-and-uncertainty}

We report micro precision, recall, and F1 for typing, taxonomy, and
non-taxonomic relations. The primary graph metric is the official exact
Graph Similarity implementation distributed with the 2026 challenge
{[}13{]}:

{[} GS=\frac{EdgeF1+NeighborhoodSimilarity+TaxonomySimilarity}{3}. {]}

The composite ((TypingF1+TaxonomyF1+GS)/3) is an internal summary and
not an official leaderboard score. Exact matching is primary. A frozen
semantic secondary analysis is reported in the supplement; its model
revision and threshold were locked before scoring.

Where reported, uncertainty intervals use 10,000 sample-cluster
bootstrap replicates. For paired system comparisons, samples are
resampled jointly within prespecified strata. We emphasize effect sizes
and support counts rather than null-hypothesis significance tests
because several slices are small and purposefully stratified.

\subsubsection{4.3 Research questions}\label{research-questions}

\begin{itemize}
\tightlist
\item
  \textbf{RQ1:} Does isolating symbolic constraints to the relation
  channel improve graph precision without changing typing or taxonomy
  output?
\item
  \textbf{RQ2:} Does guarded extraction on a controlled nonce holdout
  transfer to relation-rich official text?
\item
  \textbf{RQ3:} Do calibration and predicate routing satisfy
  prespecified independent-sample support gates?
\item
  \textbf{RQ4:} Do closed-book probeability and endpoint aliasing
  support a contamination-induced performance-inflation claim?
\end{itemize}

\begin{figure}
\centering
\pandocbounded{\includesvg[keepaspectratio]{figures/ontobridg_evidence_chain.svg}}
\caption{OntoBridg evidence chain}
\end{figure}

\textbf{Figure 1. Evaluation and decision flow.} Raw LLM generations are
processed by a deterministic relation guard and evaluated across
increasingly demanding evidence partitions. A positive controlled result
advances to independent official-text tests; failed support gates stop
expansion and deployment rather than triggering same-slice repair.

\subsection{5 Results}\label{results}

\subsubsection{5.1 Blind official-training
prefix}\label{blind-official-training-prefix}

All 120 task-system-sample units complete. Task A prompts contain
context only; Task B prompts contain only official task inputs. Table 2
shows that open generation does not improve the blind prefix. It emits
384 Task A predictions and misses Task A's sole non-taxonomic gold edge;
it emits 150 Task B predictions on a prefix containing no non-taxonomic
gold edge. Relative to closed generation, exact Graph Similarity falls
by 0.2801 on Task A and 0.2927 on Task B. The prefix therefore diagnoses
over-generation but cannot estimate relation recall reliably.

\textbf{Table 2. Gold-withheld official-training prefix (30 samples per
task).}

{\def\LTcaptype{none} % do not increment counter
\begin{longtable}[]{@{}
  >{\raggedright\arraybackslash}p{(\linewidth - 12\tabcolsep) * \real{0.1154}}
  >{\raggedright\arraybackslash}p{(\linewidth - 12\tabcolsep) * \real{0.1154}}
  >{\raggedleft\arraybackslash}p{(\linewidth - 12\tabcolsep) * \real{0.1538}}
  >{\raggedleft\arraybackslash}p{(\linewidth - 12\tabcolsep) * \real{0.1538}}
  >{\raggedleft\arraybackslash}p{(\linewidth - 12\tabcolsep) * \real{0.1538}}
  >{\raggedleft\arraybackslash}p{(\linewidth - 12\tabcolsep) * \real{0.1538}}
  >{\raggedleft\arraybackslash}p{(\linewidth - 12\tabcolsep) * \real{0.1538}}@{}}
\toprule\noalign{}
\begin{minipage}[b]{\linewidth}\raggedright
Task
\end{minipage} & \begin{minipage}[b]{\linewidth}\raggedright
System
\end{minipage} & \begin{minipage}[b]{\linewidth}\raggedleft
Typing F1
\end{minipage} & \begin{minipage}[b]{\linewidth}\raggedleft
Taxonomy F1
\end{minipage} & \begin{minipage}[b]{\linewidth}\raggedleft
NonTax F1
\end{minipage} & \begin{minipage}[b]{\linewidth}\raggedleft
Exact GS
\end{minipage} & \begin{minipage}[b]{\linewidth}\raggedleft
NonTax pred/gold
\end{minipage} \\
\midrule\noalign{}
\endhead
\bottomrule\noalign{}
\endlastfoot
A & Closed & \textbf{0.1400} & \textbf{0.4903} & 0 & \textbf{0.4404} & 0
/ 1 \\
A & Open & 0.1356 & 0.3516 & 0 & 0.1603 & 384 / 1 \\
B & Closed & 0.8824 & \textbf{0.9474} & 0 & \textbf{0.9375} & 0 / 0 \\
B & Open & \textbf{1.0000} & 0.8966 & 0 & 0.6449 & 150 / 0 \\
\end{longtable}
}

\subsubsection{5.2 Relation-channel
isolation}\label{relation-channel-isolation}

The density-stratified evaluation contains 63 Task A and 53 Task B
non-taxonomic gold edges. Table 3 reports the stable comparison: typing
and taxonomy outputs are inherited from the same open run, while only
the relation channel changes. The v1 hybrid removes 216 Task A and 42
Task B relation predictions, improving exact Graph Similarity by 0.0315
and 0.0409. V2 adds entity/predicate canonicalization and cross-sentence
evidence. It improves Task A GS by 0.0327 over open generation but is
less conservative on Task B.

No guarded hybrid recovers an exact relation on this slice. V2 generates
the one Task B relation recovered by the unguarded system at confidence
0.95, but its evidence spans are 1,604 normalized characters apart and
the frozen 600-character rule rejects it. A post-hoc wider gap accepts
that edge and additional false positives; it is reported only as a
mechanism diagnostic and does not change the primary result.

\textbf{Table 3. Relation-density-stratified hybrid comparison (15
samples per task).}

{\def\LTcaptype{none} % do not increment counter
\begin{longtable}[]{@{}
  >{\raggedright\arraybackslash}p{(\linewidth - 12\tabcolsep) * \real{0.1154}}
  >{\raggedright\arraybackslash}p{(\linewidth - 12\tabcolsep) * \real{0.1154}}
  >{\raggedleft\arraybackslash}p{(\linewidth - 12\tabcolsep) * \real{0.1538}}
  >{\raggedleft\arraybackslash}p{(\linewidth - 12\tabcolsep) * \real{0.1538}}
  >{\raggedleft\arraybackslash}p{(\linewidth - 12\tabcolsep) * \real{0.1538}}
  >{\raggedleft\arraybackslash}p{(\linewidth - 12\tabcolsep) * \real{0.1538}}
  >{\raggedleft\arraybackslash}p{(\linewidth - 12\tabcolsep) * \real{0.1538}}@{}}
\toprule\noalign{}
\begin{minipage}[b]{\linewidth}\raggedright
Task
\end{minipage} & \begin{minipage}[b]{\linewidth}\raggedright
Relation channel
\end{minipage} & \begin{minipage}[b]{\linewidth}\raggedleft
Typing F1
\end{minipage} & \begin{minipage}[b]{\linewidth}\raggedleft
Taxonomy F1
\end{minipage} & \begin{minipage}[b]{\linewidth}\raggedleft
NonTax F1
\end{minipage} & \begin{minipage}[b]{\linewidth}\raggedleft
Exact GS
\end{minipage} & \begin{minipage}[b]{\linewidth}\raggedleft
NonTax predictions
\end{minipage} \\
\midrule\noalign{}
\endhead
\bottomrule\noalign{}
\endlastfoot
A & Open & 0.3077 & 0.3475 & 0 & 0.1825 & 250 \\
A & Guarded v1 & 0.3077 & 0.3475 & 0 & \textbf{0.2140} & 34 \\
A & Guarded v2 & 0.3077 & 0.3475 & 0 & \textbf{0.2152} & 54 \\
B & Open & 0.6000 & 0.7500 & \textbf{0.0202} & 0.5598 & 46 \\
B & Guarded v1 & 0.6000 & 0.7500 & 0 & \textbf{0.6007} & 4 \\
B & Guarded v2 & 0.6000 & 0.7500 & 0 & 0.5766 & 19 \\
\end{longtable}
}

These results answer RQ1 narrowly: relation-channel isolation improves
graph precision by suppressing unmatched edges, but evidence filtering
alone does not establish relation recovery.

\subsubsection{5.3 Calibration remains
fail-closed}\label{calibration-remains-fail-closed}

The independent v2 calibration pool contains 144 eligible candidates but
only four exact positives from three samples. Calibration reduces
resubstitution Brier score from 0.8429 to 0.0270 and expected
calibration error (ECE) from 0.9024 to 0.0113, yet it fails the
five-positive-sample requirement.

Dual-model weak annotation supplies more positives but has low
relation-set agreement (micro Jaccard 0.3436), no held-out human audit,
and no overlap with the packet's single official non-taxonomic edge. In
sample-grouped five-fold evaluation, calibration reduces Brier/ECE to
0.1134/0.0045, but every fold maps the narrow 0.85-1.00 raw confidence
range to one probability block. This is base-rate shrinkage, not learned
ranking.

The untouched 24-sample v5 confirmatory slice gives the decisive test.
Base v5 yields 4 TP in 15 predictions from four samples (F1 0.1159). The
frozen calibrator improves candidate Brier/ECE from 0.6707/0.6840 to
0.2109/0.1238, but its threshold removes one TP, leaving 3 TP in 14
predictions from three samples (F1 0.0882). The confirmatory and
deployment gates therefore remain closed. The sample-cluster interval
for calibrated-minus-base F1 is {[}-0.1064, 0{]}.

\subsubsection{5.4 Contamination diagnostics do not support a global
inflation
claim}\label{contamination-diagnostics-do-not-support-a-global-inflation-claim}

The closed-book probe completes 30 samples and 214 balanced items. AUROC
is 0.9780 with a 10,000-replicate sample-cluster 95\% interval of
{[}0.9484, 0.9982{]}, showing that official facts are highly probeable.
Probeability, however, has weak and inconsistent association with
per-sample Graph Similarity.

Endpoint aliasing also has opposing effects. Original-minus-alias GS is
negative for Task A closed/open (-0.0921 and -0.1102) but positive for
Task B closed/open (+0.1557 and +0.1917). Restricting to fully
endpoint-visible samples preserves this pattern. These results show
task-dependent surface sensitivity, not verified training membership or
a general contamination-induced score increase. RQ4 is therefore
answered in the negative.

\subsubsection{5.5 Controlled extraction does not transfer to official
text}\label{controlled-extraction-does-not-transfer-to-official-text}

The nonce holdout provides the strongest positive relation result.
Guarded v2 raises exact non-taxonomic F1 from 0.2612 to 0.8571 and exact
GS from 0.1805 to 0.5690. Both systems produce zero predictions on all
60 negated controls. Guarded performance is perfect on canonical
language (F1 1.0000) but falls to 0.7119 on paraphrases, demonstrating
sensitivity even in the controlled setting.

The disjoint official-text transfer slice reverses the conclusion. It
contains 101 non-self non-taxonomic gold edges, 89 with both endpoints
visible in the compressed context. Neither system recovers one. Guarded
v2 reduces predictions from 424 to 143 and raises exact GS from 0.1946
to 0.2593, but exact relation precision, recall, and F1 remain zero. A
zero-call audit finds no exact TP among 225 raw guarded candidates and
no raw candidate using an exact gold predicate. The guard did not remove
a latent correct relation; the raw predicate family was misaligned with
the reference inventory.

\textbf{Table 4. Controlled nonce result and official-text transfer.}

{\def\LTcaptype{none} % do not increment counter
\begin{longtable}[]{@{}
  >{\raggedright\arraybackslash}p{(\linewidth - 12\tabcolsep) * \real{0.1154}}
  >{\raggedright\arraybackslash}p{(\linewidth - 12\tabcolsep) * \real{0.1154}}
  >{\raggedleft\arraybackslash}p{(\linewidth - 12\tabcolsep) * \real{0.1538}}
  >{\raggedleft\arraybackslash}p{(\linewidth - 12\tabcolsep) * \real{0.1538}}
  >{\raggedleft\arraybackslash}p{(\linewidth - 12\tabcolsep) * \real{0.1538}}
  >{\raggedleft\arraybackslash}p{(\linewidth - 12\tabcolsep) * \real{0.1538}}
  >{\raggedleft\arraybackslash}p{(\linewidth - 12\tabcolsep) * \real{0.1538}}@{}}
\toprule\noalign{}
\begin{minipage}[b]{\linewidth}\raggedright
Partition
\end{minipage} & \begin{minipage}[b]{\linewidth}\raggedright
System
\end{minipage} & \begin{minipage}[b]{\linewidth}\raggedleft
Precision
\end{minipage} & \begin{minipage}[b]{\linewidth}\raggedleft
Recall
\end{minipage} & \begin{minipage}[b]{\linewidth}\raggedleft
NonTax F1
\end{minipage} & \begin{minipage}[b]{\linewidth}\raggedleft
Exact GS
\end{minipage} & \begin{minipage}[b]{\linewidth}\raggedleft
Predictions / gold
\end{minipage} \\
\midrule\noalign{}
\endhead
\bottomrule\noalign{}
\endlastfoot
Nonce & Open & 0.2222 & 0.3167 & 0.2612 & 0.1805 & 171 / 120 \\
Nonce & Guarded v2 & \textbf{0.8644} & \textbf{0.8500} & \textbf{0.8571}
& \textbf{0.5690} & 118 / 120 \\
Official transfer & Open & 0 & 0 & 0 & 0.1946 & 424 / 101 \\
Official transfer & Guarded v2 & 0 & 0 & 0 & \textbf{0.2593} & 143 /
101 \\
\end{longtable}
}

\begin{figure}
\centering
\pandocbounded{\includesvg[keepaspectratio]{figures/ontobridg_main_results.svg}}
\caption{Controlled and official relation results}
\end{figure}

\textbf{Figure 2. Controlled success and official-text failure.} Panel A
contrasts non-taxonomic F1 on the nonce and official-transfer
partitions. Panel B reports the corresponding exact Graph Similarity.
Panel C shows that independent predicate routing produces two official
exact true positives but 138 exact false positives. Exact false
positives are scoring outcomes, not human-adjudicated semantic errors.

\subsubsection{5.6 Predicate routing gives narrow recovery and fails its
gate}\label{predicate-routing-gives-narrow-recovery-and-fails-its-gate}

The independent routing packet freezes 27 Task A samples after the
official-transfer diagnosis. It contains 69 target gold edges across
\texttt{equivalent\ class}, \texttt{disjoint\ with}, \texttt{type}, and
\texttt{database\_cross\_reference} strata and has zero sample, context,
or exact-gold overlap with prior evidence. General relation-only v3 and
routed v6 differ only in prompt identity and share the same
postprocessor.

Routed v6 raises four-stratum macro target F1 from 0 to 0.0526, but
pooled precision/recall/F1 is 0.0143/0.0290/0.0191. Both true positives
are \texttt{type} edges from two independent samples. The other three
strata remain at zero TP, so the five-positive-sample mechanism gate
fails.

\textbf{Table 5. Independent predicate-routing result.}

{\def\LTcaptype{none} % do not increment counter
\begin{longtable}[]{@{}
  >{\raggedright\arraybackslash}p{(\linewidth - 10\tabcolsep) * \real{0.1304}}
  >{\raggedleft\arraybackslash}p{(\linewidth - 10\tabcolsep) * \real{0.1739}}
  >{\raggedleft\arraybackslash}p{(\linewidth - 10\tabcolsep) * \real{0.1739}}
  >{\raggedleft\arraybackslash}p{(\linewidth - 10\tabcolsep) * \real{0.1739}}
  >{\raggedleft\arraybackslash}p{(\linewidth - 10\tabcolsep) * \real{0.1739}}
  >{\raggedleft\arraybackslash}p{(\linewidth - 10\tabcolsep) * \real{0.1739}}@{}}
\toprule\noalign{}
\begin{minipage}[b]{\linewidth}\raggedright
System
\end{minipage} & \begin{minipage}[b]{\linewidth}\raggedleft
Macro target F1
\end{minipage} & \begin{minipage}[b]{\linewidth}\raggedleft
Pooled P / R / F1
\end{minipage} & \begin{minipage}[b]{\linewidth}\raggedleft
TP / FP / FN
\end{minipage} & \begin{minipage}[b]{\linewidth}\raggedleft
TP-positive samples
\end{minipage} & \begin{minipage}[b]{\linewidth}\raggedleft
Exact GS
\end{minipage} \\
\midrule\noalign{}
\endhead
\bottomrule\noalign{}
\endlastfoot
General v3 & 0 & 0 / 0 / 0 & 0 / 0 / 69 & 0 & 0 \\
Routed v6 & \textbf{0.0526} & 0.0143 / 0.0290 / 0.0191 & 2 / 138 / 67 &
2 & 0.0016 \\
\end{longtable}
}

The post-hoc audit clarifies the failure. Of 140 accepted routed
outputs, 120 are \texttt{type}; all 138 exact false positives contain
accepted explicit or cross-sentence evidence, and 133 bind both
endpoints verbatim to the context. Routing therefore changes the failure
mode from out-of-inventory predicates to broad, evidence-bound off-gold
generation. Because official gold completeness is not established and no
independent semantic adjudication was performed, off-gold predictions
cannot be declared semantically correct or incorrect. The audit does not
reopen the gate.

\subsection{6 Discussion}\label{discussion}

\subsubsection{6.1 Evidence is necessary but not
sufficient}\label{evidence-is-necessary-but-not-sufficient}

The relation guard consistently reduces unmatched predictions and can
greatly improve graph-level scores when typing and taxonomy are held
fixed. Its success on the nonce holdout shows that the implementation
can recover unseen entity combinations from supplied text. The official
transfer and routing results show the boundary: evidence verifies that
text supports an assertion, not that the benchmark expects that
assertion or predicate serialization.

This distinction is important for ontology construction. A passage can
support many defensible type or identity statements while a reference
graph records only a subset. Under exact evaluation, these become false
positives even when both endpoints and an evidence span are valid.
Better confidence calibration cannot resolve incomplete or selectively
scoped supervision unless the calibration labels reflect that same
selection process.

\subsubsection{6.2 Why fail-closed evaluation
matters}\label{why-fail-closed-evaluation-matters}

Several intermediate analyses could have been presented optimistically:
large calibration improvements on sparse labels, a strong controlled
nonce result, a wider evidence gap that restores one relation, or a
post-hoc symmetric-edge rule. Each fails a different independence or
support requirement. The protocol turns these failures into explicit
stopping decisions. It blocks same-slice repair, prevents weak labels
from being renamed as human annotation, and avoids scaling an
intervention whose mechanism gate failed.

The resulting contribution is therefore methodological as much as
algorithmic. OntoBridg provides a reproducible sequence for separating
candidate generation, deterministic filtering, calibration, controlled
extraction, official-text transfer, and independent mechanism testing.
Negative results remain useful because their scope and stopping rules
are auditable.

\subsubsection{6.3 Implications for LLMs4OL
evaluation}\label{implications-for-llms4ol-evaluation}

First, relation-channel and core-graph metrics should be reported
separately. A model can improve typing while harming graph similarity
through relation over-generation. Second, public-training-data
experiments should be labeled as such; they do not substitute for
hidden-test or temporal evaluation. Third, contamination probes should
not create unverified clean/polluted labels. Finally, benchmarks with
sparse relation annotations may need explicit closed predicate
inventories, completeness statements, or semantic adjudication protocols
if they intend to reward plausible additional relations.

\subsection{7 Limitations, Ethics, and
Reproducibility}\label{limitations-ethics-and-reproducibility}

\textbf{Training rather than hidden test data.} All official-data
results use public 2026 training files. Gold-stratified transfer and
routing slices are mechanism diagnostics, not unbiased performance
estimates or leaderboard submissions.

\textbf{One hosted model.} The primary experiments use one
OpenAI-compatible hosted model alias. Temperature is zero and all
prompts are hashed, but provider-side weights and serving behavior are
not independently versioned. Cross-model generalization is unknown.

\textbf{Sparse and potentially incomplete gold.} Exact scoring is
appropriate for challenge comparability but cannot adjudicate plausible
off-gold relations. The routing audit demonstrates this ambiguity. No
human semantic audit is claimed.

\textbf{Synthetic holdout.} Nonce entities and relation facts are newly
generated and have no exact overlap with official data, but the relation
templates and predicates are familiar. The holdout tests controlled
new-fact extraction, not natural language coverage or unseen-predicate
discovery.

\textbf{Calibration support.} Independent exact positives remain too
sparse for deployment. Dual-model annotations are exploratory weak
labels with low agreement and no held-out human audit. They cannot open
a confirmatory gate.

\textbf{Contamination inference.} High closed-book AUROC may reflect
general public knowledge, control construction artifacts, or training
exposure. Endpoint aliasing changes lexical difficulty as well as
memorization cues. We do not identify individual training examples or
claim verified membership.

\textbf{Data licensing and privacy.} Experiments use public challenge
files and synthetic text; no personal data are introduced. The challenge
repository is MIT-licensed, while upstream ontology resources may have
additional terms. The artifact excludes third-party datasets and API
credentials. It includes code, frozen manifests, hashes, prompts,
aggregate result JSON, and zero-call replay instructions.

\textbf{Resource use.} Cached raw generations and signed postprocessing
outputs avoid repeated API calls during calibration and audit. All
reported post-hoc analyses are local. API attempt counts, retries, and
final failures are retained rather than silently discarded.

\subsection{8 Conclusion}\label{conclusion}

OntoBridg treats LLM-based ontology learning as candidate generation
followed by deterministic, fail-closed evaluation. Relation-only guards
improve graph precision and recover controlled synthetic relations, but
that success does not transfer to relation-rich official text. Predicate
routing obtains two exact \texttt{type} relations yet fails its
independent-sample support gate, while most off-gold outputs remain
evidence-bound and semantically unadjudicated. Contamination diagnostics
reveal high factual probeability without supporting a global
score-inflation claim. The defensible result is not that open-relation
ontology learning is solved, but that controlled positives, transfer
failures, and deployment decisions can be separated reproducibly. On the
current evidence, naturalistic open-relation gains remain unsupported
and every deployment gate remains closed.

\subsection{References}\label{references}

\begin{enumerate}
\def\labelenumi{\arabic{enumi}.}
\tightlist
\item
  H. Babaei Giglou, J. D'Souza, and S. Auer. LLMs4OL: Large Language
  Models for Ontology Learning. In \emph{The Semantic Web - ISWC 2023},
  pp.~408-427, 2023. https://doi.org/10.1007/978-3-031-47240-4\_22.
\item
  H. Babaei Giglou, J. D'Souza, and S. Auer. LLMs4OL 2024 Overview: The
  1st Large Language Models for Ontology Learning Challenge. \emph{Open
  Conference Proceedings}, 4:3-16, 2024.
  https://doi.org/10.52825/ocp.v4i.2473.
\item
  A. Beliaeva and T. Rahmatullaev. Heterogeneous LLM Methods for
  Ontology Learning (Few-Shot Prompting, Ensemble Typing, and
  Attention-Based Taxonomies). arXiv:2508.19428, 2025.
  https://doi.org/10.48550/arXiv.2508.19428.
\item
  S. Mishra, D. R. Meena, M. Ojha, K. P. Singh, and K. Singh. pro-team
  at LLMs4OL 2026 Tasks Flagship and Reuse: Retrieval-Augmented
  Generation and Vocabulary-Constrained Filtering for Ontology Learning.
  arXiv:2608.27101, 2026. https://doi.org/10.48550/arXiv.2608.27101.
\item
  H. Babaei Giglou, J. D'Souza, A. Aioanei, N. Mihindukulasooriya, and
  S. Auer. OntoLearner: A Modular Python Library for Ontology Learning
  with Large Language Models. arXiv:2607.01977, 2026.
\item
  M. A. Hearst. Automatic Acquisition of Hyponyms from Large Text
  Corpora. In \emph{COLING 1992}, p.~539, 1992.
  https://doi.org/10.3115/992133.992154.
\item
  E. Jimenez-Ruiz and B. Cuenca Grau. LogMap: Logic-Based and Scalable
  Ontology Matching. In \emph{The Semantic Web - ISWC 2011},
  pp.~273-288, 2011. https://doi.org/10.1007/978-3-642-25073-6\_18.
\item
  D. Faria, C. Pesquita, E. Santos, M. Palmonari, I. F. Cruz, and F. M.
  Couto. The AgreementMakerLight Ontology Matching System. In \emph{OTM
  2013 Conferences}, pp.~527-541, 2013.
  https://doi.org/10.1007/978-3-642-41030-7\_38.
\item
  Y. He, J. Chen, D. Antonyrajah, and I. Horrocks. BERTMap: A BERT-Based
  Ontology Alignment System. \emph{Proceedings of the AAAI Conference on
  Artificial Intelligence}, 36(5):5684-5691, 2022.
  https://doi.org/10.1609/aaai.v36i5.20510.
\item
  Y. Oren, N. Meister, N. Chatterji, F. Ladhak, and T. B. Hashimoto.
  Proving Test Set Contamination in Black-Box Language Models.
  arXiv:2310.17623, 2023. https://doi.org/10.48550/arXiv.2310.17623.
\item
  S. Golchin and M. Surdeanu. Time Travel in LLMs: Tracing Data
  Contamination in Large Language Models. In \emph{ICLR 2024}, 2024.
  arXiv:2308.08493.
\item
  C. Li and J. Flanigan. Task Contamination: Language Models May Not Be
  Few-Shot Anymore. \emph{Proceedings of the AAAI Conference on
  Artificial Intelligence}, 38(16):18471-18480, 2024.
  https://doi.org/10.1609/aaai.v38i16.29808.
\item
  LLMs4OL Organizers. LLMs4OL Challenge Repository: 2026 Tasks and
  Metrics. GitHub repository, 2026.
  https://github.com/LLMs4OL/LLMs4OL-Challenge (accessed 2026-09-28).
\end{enumerate}

\end{document}
